\section{Which intersection fixes the exponent}
\label{sec:h-rhoi}

Hold a comoving population of surfaces inside an already expanding
isotropic volume, as a kinematic comparison.
Area density scales as \(a^{-1}\).
Length density of pairwise intersection lines scales as \(a^{-2}\).
Number density of triple points scales as \(a^{-3}\).
The relation \(a\propto\rho_I^{-1/3}\) is the triple-point count.
Line density would give \(a\propto\rho_L^{-1/2}\).
Sheet area would give \(a\propto\rho_S^{-1}\).

The stitch intensity \(\lambda\) is a number of stitches per node.
That is a point process on the sheet, not a triple point in the
three-volume.
On the chain the derived ruler is \(\ell_W\propto\lambda^{-1/2}\)
\eqref{eq:ell}, so the transverse scale tracks the exponent \(-1/2\).
Replacing \(-1/2\) by \(-1/3\) changes \(a_\perp(\lambda)\) by a power.
You could keep the Poisson exponent unless a different intersection is
the one being counted, and name that intersection when the exponent
changes.
Section~\ref{sec:expand} is the same separation, written without a
preferred density.
